% Section 2 - Classes
% Roberto Masocco <roberto.masocco@uniroma2.it>
% April 12, 2026

% ### Classes ###
\section{Classes}
\graphicspath{{figs/section2/}}

% --- Classes ---
\begin{frame}{Classes}{\texttt{class} vs \texttt{struct}}
	\justifying
	C++ keeps the \texttt{struct} keyword from C, but gives it new powers: \textbg{member functions}, \textbg{constructors}, and \textbg{access control}.\\
	\bigskip
	It also introduces \texttt{class}. The \textbg{only difference} between the two is the \textbg{default access specifier}:
	\begin{itemize}
		\item members of a \texttt{struct} are \textbg{public} by default;
		\item members of a \texttt{class} are \textbg{private} by default.
	\end{itemize}
	\bigskip
	By convention:
	\begin{itemize}
		\item \texttt{struct} is used for \textbg{plain data aggregates} with no significant behavior;
		\item \texttt{class} is used when \textbg{encapsulation} and \textbg{operations} matter.
	\end{itemize}
\end{frame}
\begin{frame}[fragile]{Classes}
  \begin{columns}
    \column{.9\linewidth}
		\begin{lstlisting}[style=codebeamer, language=C++, caption=Example of a C++ class.]
class MyClass : public ParentClass
{
public:
  MyClass();
  ~MyClass();
  int get_int() { return variable_int_member_; }
  // ...
protected:
  // ...
private:
  int variable_int_member_ = 0;
  double variable_double_member_ = 0.0;
  // ...
};
    \end{lstlisting}
  \end{columns}
\end{frame}
\begin{frame}{Classes}{Access specifiers}
	\justifying
	Access specifiers control the \textbg{visibility} of class members to the outside world.\\
	\bigskip
	There are three:
	\begin{itemize}
		\item \texttt{public}: accessible from anywhere;
		\item \texttt{protected}: accessible from the class itself and from \textbg{derived classes};
		\item \texttt{private}: accessible \textbg{only} from the class itself, \textbg{not even from derived classes}.
	\end{itemize}
	\bigskip
	A specifier applies to all members declared \textbg{below it}, until the next specifier.\\
	\bigskip
	Good practice is to expose only what is \textbg{necessary}:
	\begin{itemize}
		\item \textbg{data members} are almost always \texttt{private}, accessible via \textbg{getters}/\textbg{setters};
		\item the \textbg{public interface} is kept minimal and stable.
	\end{itemize}
\end{frame}
\begin{frame}{Classes}{Member functions and the \texttt{this} pointer}
	\justifying
	Functions declared inside a class are called \textbg{member functions} (or \emph{methods}).\\
  They have \textbg{direct access to all data members of the class}, regardless of access specifiers.\\
	\bigskip
	Inside any (non-static) member function, the keyword \texttt{this} is an implicit \textbg{pointer to the object on which the function was called}.\\
	\bigskip
	Member functions that \textbg{do not modify the object} should be declared \texttt{const}: the \textbg{compiler} then enforces the constraint.
\end{frame}
\begin{frame}{Classes}{Constructors and destructors}
	\justifying
	Every class can define \textbg{special member functions} that the compiler calls automatically.\\
	\bigskip
	The \textbg{constructor} is called when an object is \textbg{created}:
	\begin{itemize}
		\item its name is the class name, it has no return type;
		\item it initializes data members and acquires resources;
		\item a class may have \textbg{multiple constructors} with different parameter lists (\textbg{overloading});
		\item if none is defined, the compiler generates a \textbg{default constructor}.
	\end{itemize}
	\bigskip
	The \textbg{destructor} is called when an object is \textbg{destroyed}:
	\begin{itemize}
		\item its name is the class name preceded by \texttt{\textasciitilde};
		\item it \textbg{releases resources} and \textbg{performs cleanup};
		\item there is \textbg{exactly one} destructor per class, it takes no arguments.
	\end{itemize}
\end{frame}
\begin{frame}{Classes}{Constructors and destructors}
	\begin{alertblock}{Guarantee}
		\justifying
		\textbf{The destructor is always called}, regardless of how the object's scope ends — even in the presence of exceptions.
	\end{alertblock}
\end{frame}
\begin{frame}{Classes}{Separating declaration from definition}
	\justifying
	In C++, a class is typically \textbg{declared} in a header file (\texttt{.hpp}) and \textbg{defined} in a source file (\texttt{.cpp}).\\
	\bigskip
	The \textbg{header} contains
	\begin{itemize}
		\item the \textbg{class declaration} with all member names and types;
		\item \texttt{inline} \textbg{function definitions} (small, frequently called functions).
	\end{itemize}
	\bigskip
	The \textbg{source file} contains
	\begin{itemize}
		\item the full \textbg{definitions} of all non-inline member functions;
		\item member function names are \textbg{prefixed} with \texttt{ClassName::} (their \textbg{namespace}).
	\end{itemize}
	\bigskip
	\begin{block}{Why split?}
		\justifying
		The header is the \textbf{public contract} of the class — it is what other translation units include.
		Keeping implementation in \texttt{.cpp} files reduces compile-time dependencies and \textbf{hides internals}.
	\end{block}
\end{frame}
